\documentclass[../main.tex]{subfiles}

\begin{document}

\justifying

\section{Pruebas del gel con el potencial de 3 cuerpos bien tabulado}

En el directorio\ldots se hicieron las pruebas y se comprobó que el potencial de 3 cuerpos está bien tabulado.
Solo hay un error numérico en la fuerza que hace que qued euna fuerza residual \emph{despreciable} al compararlo con la magnitud de la fuerza de enlace\footnote{No he hecho na comparación cuantitativa y no creo que sea necesario hacerla}.

Ahora se analizará el código del gel polimérico.

\subsection{Directory: }

\newpage

\section{Los patches más adentro}

Revisando el código\footnote{Había varios errores en la definición de los potenciales.
Se usaba el mismo parametro de bond entre patches, pero la distancias no correspondian.
}, los patches de los monomeros se encontraban en la superficie del cascarón de la partícula central.
\(r_{\mathrm{PB}} = (\pm r_{\mathrm{CP}}, 0, 0)\).
Mientrás que las coordenadas de los patches en el CL están dados por:
\begin{align*}
    r_{\mathrm{PA}1} &= (0,0,0.46) \\
    r_{\mathrm{PA}2} &= (-0.22,-0.38,-0.15) \\
    r_{\mathrm{PA}3} &= (-0.22,0.38,-0.15) \\
    r_{\mathrm{PA}4} &= (0.43,0,-0.15)
\end{align*}
Esas coordenadas coinciden con las de una esfera de radio \(0.46\) circunscrita en un tetraedro.

Se realizarán los siguiente cambios.
Los patches estarána una distancia \(D_p\) con respecto al centro de la partícula central.
Es decir que los patches de los monomeros estarán ubicados en \(r_{\mathrm{PB}} = (\pm D_p, 0, 0)\).
Mientrás que los patches del CL, en coordenadas esféricas, estarán en,
\begin{align*}
    r_{\mathrm{PA}1} &= (D_p,0,0) \\
    r_{\mathrm{PA}2} &= (D_p,\acos(-1/3),0) \\
    r_{\mathrm{PA}3} &= (D_p,\acos(-1/3),2\pi/3) \\
    r_{\mathrm{PA}4} &= (D_p,\acos(-1/3),4\pi/3)
\end{align*}
En coordenadas cartesianas es,
\begin{align*}
    r_{\mathrm{PA}1} &= (0,0,D_p) \\
    r_{\mathrm{PA}2} &= (\frac{2\sqrt{2}}{3}D_p,0,-\frac{D_p}{3}) \\
    r_{\mathrm{PA}3} &= (-\frac{\sqrt{2}}{3}R,\frac{\sqrt{6}}{3}D_p,\frac{D_p}{3}) \\
    r_{\mathrm{PA}4} &= (-\frac{\sqrt{2}}{3}R,-\frac{\sqrt{6}}{3}D_p,\frac{D_p}{3})
\end{align*}
Al sustituir los valores numéricos y convirtiendo a coordenadas cartesianas,
\begin{align*}
    r_{\mathrm{PA}1} &= (0,0,0.4) \\
    r_{\mathrm{PA}2} &= (0.377,0,-0.133) \\
    r_{\mathrm{PA}3} &= (-0.188,0.326,-0.133) \\
    r_{\mathrm{PA}4} &= (-0.188,-0.326,-0.133)
\end{align*}



\subsection{Directory: }



\end{document}
